\subsection{Sanitizeri u Clang-u}\label{clang-sanitizeri}

Sanitizeri su instrumentacije koje kompilator ubacuje u program. Program se zatim pokreće kao obična izvršna datoteka, ali dodatni kod proverava memorijske pristupe, aritmetiku, curenje memorije ili pristup deljenim podacima. Zbog toga su sanitizeri pogodni za lokalni razvoj i CI provere. U odnosu na Valgrind, zahtevaju posebno prevođenje programa, ali se često izvršavaju brže i daju precizan stek poziva za konkretno izvršavanje.

Primeri u ovom direktorijumu koriste Clang:

\begin{minted}{bash}
$ make
\end{minted}

\paragraph{ASan}\label{addresssanitizer}

\href{https://clang.llvm.org/docs/AddressSanitizer.html}{AddressSanitizer} (ASan) detektuje neispravne pristupe memoriji kao što su prekoračenje hip bloka, prekoračenje steka, upotreba memorije posle oslobađanja i neispravno oslobađanje memorije. Primer ispod alocira prostor za string bez završnog nul karaktera, a zatim kopira i taj bajt.

\codefile{01_asan_heap_overflow.c}{c}

\codefile{01_asan_heap_overflow_fixed.c}{c}

\begin{minted}{bash}
$ ./asan_bad.out
$ ./asan_fixed.out
\end{minted}

U izveštaju loše verzije tražiti red \texttt{heap-buffer-overflow}. Stek poziva pokazuje gde je upis van granica nastao, a deo izveštaja sa opisom alociranog regiona pokazuje koliko bajtova je zaista rezervisano.

\paragraph{UBSan}\label{undefinedbehaviorsanitizer}

\href{https://clang.llvm.org/docs/UndefinedBehaviorSanitizer.html}{UndefinedBehaviorSanitizer} (UBSan) prijavljuje ponašanje koje C i C++ standard ne definišu. To obuhvata potpisano celobrojno prekoračenje, neispravne pomeraje bitova, dereferenciranje neporavnatih pokazivača i slične greške. U ovom primeru množenje prelazi opseg tipa \texttt{int}.

\codefile{02_ubsan_signed_overflow.c}{c}

\codefile{02_ubsan_signed_overflow_fixed.c}{c}

\begin{minted}{bash}
$ ./ubsan_bad.out
$ ./ubsan_fixed.out
\end{minted}

Makefile koristi opciju \texttt{-fno-sanitize-recover=undefined} da se program zaustavi posle prijave. U CI proverama je to korisno jer greška dobija nenulti izlazni status.

\paragraph{LSan}\label{leaksanitizer}

\href{https://clang.llvm.org/docs/LeakSanitizer.html}{LeakSanitizer} (LSan) proverava curenje memorije na kraju procesa. Može da se koristi samostalno opcijom \texttt{-fsanitize=leak}, a često se koristi i zajedno sa ASan instrumentacijom. Loša verzija učitava tri linije i gubi pokazivače na alocirane kopije.

\codefile{03_lsan_leak.c}{c}

\codefile{03_lsan_leak_fixed.c}{c}

\begin{minted}{bash}
$ ./lsan_bad.out
$ ./lsan_fixed.out
\end{minted}

Kod loše verzije očekuje se izveštaj sa tekstom \texttt{detected memory leaks}. Popravljena verzija oslobađa svaku kopiju pre izlaska iz funkcije.

\paragraph{TSan}\label{threadsanitizer}

\href{https://clang.llvm.org/docs/ThreadSanitizer.html}{ThreadSanitizer} (TSan) detektuje trke za podacima u programima sa nitima. Primer ima dve niti koje uvećavaju istu globalnu promenljivu bez sinhronizacije.

\codefile{04_tsan_data_race.c}{c}

\codefile{04_tsan_data_race_fixed.c}{c}

\begin{minted}{bash}
$ ./tsan_bad.out
$ ./tsan_fixed.out
\end{minted}

U lošoj verziji očekuje se \texttt{data race} izveštaj. Popravljena verzija štiti deljeni brojač muteksom, pa svaki pristup promenljivoj \texttt{counter} ima istu sinhronizacionu disciplinu.

\paragraph{MemorySanitizer kao napredna vežba}\label{memorysanitizer-napredno}

\href{https://clang.llvm.org/docs/MemorySanitizer.html}{MemorySanitizer} (MSan) detektuje čitanje neinicijalizovane memorije. U praksi često zahteva da su i biblioteke koje program koristi prevedene sa MSan instrumentacijom. Zbog toga ga ne koristimo kao osnovnu lokalnu vežbu u ovom poglavlju. Kao naprednu vežbu, izolovati mali program bez spoljašnjih zavisnosti i prevesti ga komandom:

\begin{minted}{bash}
$ clang -g -O1 -fsanitize=memory -fno-omit-frame-pointer example.c -o example
\end{minted}

\paragraph{Mesta za proširenje}\label{clang-prosirenja-backlog}

Sledeće teme su dobar nastavak nakon Clang sanitizer vežbi:

\begin{enumerate}
\tightlist
\item
  \textbf{Kombinovanje sanitizera i fuzz testiranja}: ASan i UBSan se često uključuju tokom fuzz testiranja kako bi se padovi programa pretvorili u precizne izveštaje.
\item
  \textbf{LLVM pokrivenost izvornog koda}: \texttt{llvm-profdata} i \texttt{llvm-cov} mogu povezati izvršene testove sa pokrivenošću izvornog koda.
\item
  \textbf{Podešavanje}: opcije kao što su \texttt{ASAN\_OPTIONS} i \texttt{TSAN\_OPTIONS} kontrolišu izveštaje, supresije i ponašanje pri grešci.
\item
  \textbf{CI konfiguracija}: sanitizer build treba držati odvojen od produkcionog build-a i pokretati ga nad testovima koji imaju dovoljno pokrivenosti.
\end{enumerate}
